\section{Learning the marginal continuation of the original Bellman operator}
\label{sec:marginal121}

The preceding comparison methods certify an acquired actor.  We now change
what the critic learns.  The target remains the original two-date
investment problem in Section~\ref{sec:continuous95}, with its continuous
innovation law, capacity, discount and nested AVaR preference.  The
construction first evaluates the original risk operator and then learns
its action marginal.  In particular, it does not replace the optimizing
risk density by the density selected by a fitted terminal continuation.

Write \(a=A(x)v\), \(v\in[0,1]\), and define
\begin{align}
 R_\alpha(x,v)
   &=\beta\rho_\alpha\{V_1(F(x,A(x)v,Z))\},\label{eq:rootcritic121}\\
 Q_x(A(x)v)&=S(x)+h(A(x)v)+R_\alpha(x,v),\nonumber\\
 H_\alpha(x,v)&=A(x)^{-1}\partial_vR_\alpha(x,v).\nonumber
\end{align}
Thus \(H_\alpha\) is a marginal continuation cost per unit of the
original action.  The known stage cost is retained analytically.
The original optimal last-date value \(V_1=T_1g\) is the target of the
paid acquisition calculation.  Deployment still uses the same measurable
last-date controller as the retained study.  Its uniform guarantee gives
\begin{equation}\label{eq:future121}
 J_x(a)=Q_x(a)+\eta_x(a),\qquad
 0\leq\eta_x(a)\leq\delta:=\beta\,2^{-20}.
\end{equation}
No differentiability or convexity of \(\eta_x\) is assumed.

Derivative-based economic dynamic programming, derivative supervision
and decision-aware learning have established precedents
\citep{CaiJudd2015,Czarnecki2017Sobolev,Donti2017Task,ElmachtoubGrigas2022}.
The action-null quotient principle is also retained from R52, and the
constrained directional regret bound from Theorem~\ref{thm:action99}.
R67 already supplies a certified fixed-feature readout and a dual gap.
The additional construction here supplies regularity constants for the
original recursive AVaR target, outward derivative labels made from paid
finite-risk values, a finite-class training certificate, and a
feature-aware extension from those labels to off-design decisions.
These ingredients apply to conventional and neural representations with
the same information and feature restrictions.

\subsection{Regularity of the original integrated risk marginal}

Throughout this section \(d=2\), \(R=1/32\), \(\beta=15/16\),
\(\alpha\in\{1/2,3/4\}\), and \(t=1-\alpha\).  Retain the original
state-gradient constant from Theorem~\ref{thm:modulus120}:
\begin{align}
 M_\alpha&=\frac{1413}{128}
   +\beta\left\{\frac{127}{16}
               +\frac{19625\alpha}{20480t}\right\},\nonumber\\
 m_z&=\frac{181}{64},\qquad C_\alpha=\frac45M_\alpha,\nonumber\\
 D_\alpha&=\frac{5M_\alpha}{16}
                  +\frac{C_\alpha^2\alpha}{m_zt},\qquad
 L_\alpha=5+\beta D_\alpha.\label{eq:marginalconstants121}
\end{align}
All these constants are derived from the original primitives.
They are not fitted from the acquisition or evaluation sample.

\begin{theorem}[An original-risk marginal with explicit boundary control]
\label{thm:rootmarginal121}
The function \(H_\alpha\) in \eqref{eq:rootcritic121} is continuous on
\(\Omega=[0,1]^2\times[0,1]\), with the inward derivative convention at
the action endpoints.  For each state, \(Q_x\) is continuously
differentiable and \(2\)-strongly convex on \([0,A(x)]\), and
\[
 |Q_x'(a)-Q_x'(b)|\leq L_\alpha|a-b|.
\]
Consequently its normalized action curvature lies between
\(2A(x)^2\) and \(L_\alpha A(x)^2\).
For \(u=(x,v)\) and \(u'=(x',v')\),
\begin{align}
 |H_\alpha(u)-H_\alpha(u')|
 &\leq \Lambda_\alpha\|u-u'\|_2,\label{eq:rootmarginalmodulus121}\\
 \Lambda_\alpha
 &=\beta\left\{\frac{441M_\alpha}{1024}
       +\frac{147M_\alpha^2\alpha}{160m_zt}\right\}.\nonumber
\end{align}
These statements include a minimizing central risk interval at either
endpoint of its feasible set and a binding economic capacity.
\end{theorem}

The proof uses the original \(M_\alpha\)-Lipschitz gradient of \(V_1\)
and the \(m_z\)-strong convexity of
\(z\mapsto V_1(F(x,a,z))\).  A central interval of length \(2R\alpha\)
has a unique minimizing location.  Its constrained first-order
inequalities, rather than an assumed interior root, control changes
in the adverse tail.  That step is essential: controlling a fitted
terminal continuation in one action direction alone need not control
which risks receive AVaR weight.

\subsection{Paid derivative labels and finite-class training}

At a declared acquisition site \((x,v)\), set \(a=A(x)v\).
Suppose an original-\(Q_x\) oracle returns outward value intervals at
\(a\) and the available inward probes \(a\pm d\), where \(d>0\)
is fixed before their outcomes and each queried action is feasible.
Write \(\mathcal I_0,\mathcal I_+,\mathcal I_-\) for these intervals,
and use outward interval subtraction and division to form
\[
 \mathcal S_+=(\mathcal I_+-\mathcal I_0)/d,\qquad
 \mathcal S_-=(\mathcal I_0-\mathcal I_-)/d,\qquad
 \mathcal S_c=(\mathcal I_+-\mathcal I_-)/(2d).
\]

\begin{proposition}[Finite-risk secants enclose the original marginal]
\label{prop:marginalprobe121}
Every available interval in the following list contains \(Q_x'(a)\):
\begin{align}
 \mathcal G_+
   &=\mathcal S_++[-L_\alpha d/2,-d],\nonumber\\
 \mathcal G_-
   &=\mathcal S_-+[d,L_\alpha d/2],\label{eq:secantlabels121}\\
 \mathcal G_c
   &=\mathcal S_c+
        [-(L_\alpha-2)d/4,(L_\alpha-2)d/4].\nonumber
\end{align}
The central interval requires both probes.  The forward interval alone
is valid at \(a=0\), and the backward interval alone at \(a=A(x)\).
Their outward intersection, followed by subtraction of the known
\(h'(a)\), gives a valid label
\([\ell_i,u_i]\) for \(H_\alpha(x,v)\).
\end{proposition}

The experiment uses the prospectively specified normalized probe
\(2^{-12}\), hence \(d=A(x)2^{-12}\).  Its five main action sites
are \(v\in\{0,1/4,1/2,3/4,1\}\).  The three interior sites buy
three original-\(Q\) values each, and the endpoints buy two each:
there are thirteen paid contexts at each state.  The nested state
meshes have \(5^2\) and \(9^2\) sites, giving 325 and 1053
contexts per risk level and standalone acquisition budget.
The protocol specifies how each continuous-risk interval is
constructed and charges all optimal-continuation solves, probes and
enclosure work.  The allowance \(\delta\) in \eqref{eq:future121}
does not enter an optimal-\(Q\) derivative label; it enters the
subsequent implemented-policy statement once.

For a declared family of continuous piecewise-polynomial features,
use the anchored critic
\begin{equation}\label{eq:anchoredcritic121}
 \widehat R_\theta(x,v)=\widehat R_0(x)
       +\sum_{j=1}^p\theta_j
                   \{\phi_j(x,v)-\phi_j(x,0)\}.
\end{equation}
The implementation stores an approximation \(\widehat Q_0(x)\)
to \(Q_x(0)\).  In \eqref{eq:anchoredcritic121} this means
\(\widehat R_0(x)=\widehat Q_0(x)-S(x)\), so the full action
critic is \(\widehat Q_0(x)+h(A(x)v)\) plus the learned
anchored difference.  The state-only anchor affects a displayed
value level and cancels from every current action comparison.
At a label site, let
\(A_{ij}=A(x_i)^{-1}D_v\phi_j(x_i,v_i)\).
Include both feasible one-sided rows at an interior kink and only
the inward row at \(v=0\) or \(v=1\).
The same label is attached to each row.  Let the coefficient box
\(\Theta=\prod_j[\underline\theta_j,\overline\theta_j]\) be fixed
before acquisition.  The action-marginal training objective is
\begin{equation}\label{eq:robustlp121}
 T(\theta)=\max_i\{A_i\theta-\ell_i,\ u_i-A_i\theta\},
 \qquad T_*=\min_{\theta\in\Theta}T(\theta).
\end{equation}
It is the largest possible node error over the retained outward
labels, in original-action units.  The value-loss control uses the
same box, features and complete paid bundle, and fits the anchored
values by bounded least squares.  A smaller value loss is not
presumed to minimize \eqref{eq:robustlp121}.

\begin{proposition}[A checkable global training gap in the declared box]
\label{prop:boxtraining121}
The minimum \(T_*\) exists.  Any exactly checked
\(\widehat\theta\in\Theta\) gives the primal upper bound
\(P=T(\widehat\theta)\).  Let \(y_i^+,y_i^-\geq0\) be rational
weights whose combined sum is one, and set
\[
 r=\sum_i(y_i^+-y_i^-)A_i,\qquad
 B_y=-\sum_i\ell_i y_i^++\sum_i u_i y_i^-
       +\sum_j\min\{\underline\theta_jr_j,
                         \overline\theta_jr_j\}.
 \]
Then
\begin{equation}\label{eq:traininggap121}
 \max\{0,B_y\}\leq T_*\leq P,\qquad
 0\leq P-T_*\leq P-\max\{0,B_y\}.
\end{equation}
In particular, no exact dual stationarity equations or solver
status are needed for the lower bound.  Exact verification must
bind the actual stored coefficients, feature rows, labels and box.
\end{proposition}

The same independent rational audit checks the value-loss
control's optimization gap by minimizing the affine gradient
support of its convex sum of squares over the declared box.
Proposition~\ref{prop:valueloss121} supplies that lower bound
and its proof.  Thus neither training objective is declared
globally solved from a floating optimizer's success flag.

The declared box is \([-128,128]\) except for the nonnegative
quadratic curvature coefficient and the nonnegative nonlinear
ReLU or matched-hinge coefficients.  A signed coordinate contributes
\(-128|r_j|\) to \(B_y\); a nonnegative one contributes
\(\min\{0,128r_j\}\).
The quadratic, polynomial, spline, tensor-grid and matched-hinge
controls face the same acquisition interface.  A coordinate-loop
hinge carrying the same features and coefficients represents the
same function as its ReLU encoding.  This experiment can identify
the consequence of a learning loss or encoding cost, not a
neural-specific approximation advantage from that equality.

\subsection{From finite labels to an off-design policy certificate}

A node training certificate does not supply uniform derivative
accuracy.  To make the remaining calculation explicit, take a
finite cover \(\mathcal C\) of \(\Omega\).  For each cell \(C\),
construct intervals
\[
 d_j(C)\supseteq
  \{A(x)^{-1}D_v\phi_j(x,v):(x,v)\in C\},
 \]
including every relevant one-sided value.  Let
\([\underline h_\theta(C),\overline h_\theta(C)]\) enclose
\(\sum_j\theta_jd_j(C)\) by interval arithmetic, and let
\(r_i(C)\geq\sup_{u\in C}\|u-\xi_i\|_2\), where
\(\xi_i=(x_i,v_i)\) is the parameter node.  Define
\begin{align}
 \underline H(C)&=\max_i\{\ell_i-\Lambda_\alpha r_i(C)\},\nonumber\\
 \overline H(C)&=\min_i\{u_i+\Lambda_\alpha r_i(C)\},\nonumber\\
 \epsilon(C)&=\max\{0,\overline h_\theta(C)-\underline H(C),
                     \overline H(C)-\underline h_\theta(C)\}.
 \label{eq:finitecover121}
\end{align}
Using only a declared subset of label sites in the maximum and
minimum remains valid.  Empty target intersections indicate a
failed numerical enclosure and do not authorize selection of a
favorable endpoint.

\begin{theorem}[Feature-aware continuation and decision control]
\label{thm:quotientlearning121}
Suppose the label and feature-range inequalities above are valid.
For each cell \(C\), the original-action marginal error of
\eqref{eq:anchoredcritic121} is at most \(\epsilon(C)\), including
one-sided limits.  Fix a state \(x\), and let
\(\epsilon_x(v)\) be any measurable choice of a valid cell bound
at \((x,v)\).  For feasible original actions \(a,b\), put
\begin{align}
 \mathcal D_x(a,b)
   &=A(x)\int_{\min(a,b)/A(x)}^{\max(a,b)/A(x)}
                          \epsilon_x(v)\,dv,\nonumber\\
 \widehat q_x(a)&=S(x)+h(a)+
        \widehat R_\theta(x,a/A(x)).\label{eq:pathcertificate121}
\end{align}
Then
\begin{align}
 &\left|Q_x(b)-Q_x(a)
          -\{\widehat q_x(b)-\widehat q_x(a)\}\right|
       \leq\mathcal D_x(a,b),\label{eq:quotientcontrast121}\\
 &J_x(b)-J_x(a)\leq
       \widehat q_x(b)-\widehat q_x(a)+\mathcal D_x(a,b)+\delta,
       \nonumber\\
 &0\leq J_x(a)-V_0(x)\leq \delta+
       \sup_{0\leq b\leq A(x)}
       \{\widehat q_x(a)-\widehat q_x(b)+\mathcal D_x(a,b)\}.
       \label{eq:learnedallaction121}
\end{align}
Any certified upper bound on the final supremum remains valid.
These conclusions do not require a convex fitted critic.
\end{theorem}

The finite cover pays for the actual feature geometry.  An anchored
ReLU feature has derivative zero or its declared action coefficient;
both sides enter a crossing cell.  Polynomial and spline pieces
have computable polynomial derivative ranges.  Tensor hats have
one-sided action derivatives at their knots.  Dividing their
numerator bounds by the positive capacity interval is valid because
\(A(x)\geq1/8\).  The supplement gives constructive range rules
for every declared family.  Consequently
\eqref{eq:finitecover121} does not infer a smooth fitted derivative
from smoothness of the original target.

The bound can be conservative at the fixed training meshes.
Unresolved cells, paid refinements and fallback original-risk
qualification remain part of the numerical record.  No off-design
certificate is credited merely because a finite-dimensional fit
has a small node loss.

\begin{proposition}[Capacity-sensitive actor extraction]
\label{prop:marginalpolicy121}
Suppose a feasible returned action \(a\) has an outward original
derivative interval \([g_-,g_+]\).  Define
\begin{align}
 \mathcal K_x(a;g)&=
   -\min_{d\in[-a,A(x)-a]}\{gd+d^2\},\nonumber\\
 K_x(a)&=\max\{\mathcal K_x(a;g_-),
                         \mathcal K_x(a;g_+)\}.
 \label{eq:marginalkkt121}
\end{align}
Then
\[
 Q_x(a)-V_0(x)\leq K_x(a),\qquad
 J_x(a)-V_0(x)\leq K_x(a)+\delta,\qquad
 |a-a_0^*(x)|^2\leq K_x(a).
\]
The scalar minimum is attained at
\(\operatorname{proj}_{[-a,A(x)-a]}(-g/2)\).
At \(a=0\), \(g_-\geq0\) implies \(K_x(a)=0\);
at \(a=A(x)\), \(g_+\leq0\) implies \(K_x(a)=0\).
If the fitted actor has a checked constrained stationarity
residual at most \(r\), and every required one-sided marginal
error is at most \(\epsilon\), the interval supplied by that
error bound can be chosen so that
\[
 K_x(a)\leq (r+\epsilon)^2/4.
\]
\end{proposition}

This is the retained directional strong-convexity bound applied
to a newly acquired and extended marginal label.  It makes a
capacity action with the correct marginal sign economical to
certify.  A numerical polynomial root or an optimizer status
alone does not supply the stated residual; the returned action,
piecewise derivatives and rounding must be checked or covered
by the original certificate service.  Similarly, neither a
training-gap certificate nor a certified action identifies a
computational adoption region without the full acquisition and
decision-service costs and the common-information comparisons.

\subsection{The anchor that a recursive Bellman operator must retain}

An action-independent level cancels at its current state.
It need not cancel when that learned value is exported as
the next continuation.  The original investment transition
gives an exact example.  If
\(W_1(y)=V_1(y)+\gamma(y_1+y_2)\), the zero-sum shock and
cash invariance imply
\[
 Q_x^{W_1}(a)-Q_x(a)
    =\beta\gamma\{F_1(x,a,0)+F_2(x,a,0)\}.
\]
The added root marginal is \(3\beta\gamma/4\).
The same state-dependent level could leave a last-date
action minimization unchanged.  Thus discarding that
anchor would change the recursive decision problem.

The complete-value statement in
Theorem~\ref{thm:recursiveanchor121} supplies the required
account.  Let \(W_{t+1}\) be an exported continuation
with signed uniform errors
\(E_{t+1}^-\leq W_{t+1}-V_{t+1}\leq E_{t+1}^+\).
Suppose the acquired action critic is compared with
the original Bellman objective using that \(W_{t+1}\),
and its retained anchor and action-contrast errors
give full-level error bounds
\(d_t^-\leq\widehat Q_t-Q_t^{W_{t+1}}\leq d_t^+\).
If the exported value has an additional checked
minimization error at most \(\zeta_t\) in absolute
value, then a monetary recursive risk map gives
\begin{align}
 E_t^-&=d_t^-+\beta_tE_{t+1}^- -\zeta_t,\nonumber\\
 E_t^+&=d_t^++\beta_tE_{t+1}^+ +\zeta_t.
 \label{eq:anchorrecursion121}
\end{align}
A local certified action gap \(g_t\) under
\(Q_t^{W_{t+1}}\) and future policy bound \(P_{t+1}\)
give
\[
 P_t\leq g_t+\beta_t(E_{t+1}^+-E_{t+1}^-)
                         +\beta_tP_{t+1}.
\]
In the present paid two-date construction the root
target uses the original \(V_1\), so its continuation
approximation term is zero; the unchanged terminal
controller contributes \(\beta\varepsilon_1=\delta\).
The supplementary theorem separately states the
bound for a noncash-invariant recursive preference
under an explicit continuation Lipschitz assumption.
Such a preference does not inherit the monetary
shift identity by notation.
